\section{Introduction}
\label{sec:intro}

Short-term network-wide traffic forecasting underpins signal control, ramp metering, incident management and traveller information. Deep spatio-temporal models such as diffusion-convolutional recurrent networks \citep{li2018dcrnn}, spatio-temporal graph convolutional networks \citep{yu2018stgcn,wu2019graphwavenet,bai2020agcrn} and, more recently, spatio-temporal Transformers \citep{vaswani2017attention,stlaformer} have substantially improved accuracy on benchmark data. These models are almost always trained and evaluated on histories in which every detector reports every interval. In operation this assumption is routinely violated: loop detectors fail individually, communication links drop whole corridors, maintenance removes sensors for hours, and power or data-pipeline incidents remove entire regions or one travel direction at once \citep{cui2020missing,chen2019bayesian}.

Two families of remedies exist. \emph{Impute-then-forecast} pipelines first reconstruct the missing history with a dedicated model \citep{cao2018brits,cini2022grin,tashiro2021csdi,chen2019bayesian} and then run a standard forecaster. \emph{Direct} approaches feed the incomplete history together with its mask to a forecaster that has been made missing-aware \citep{che2018grud,tang2020joint,cui2020missing}. Both families treat missingness as a loss of input information that has to be compensated inside the network. What they do not model explicitly is the consequence of missingness on the \emph{output} of a given forecaster. For a fixed model, removing inputs displaces the forecast in a systematic way that depends on which entries are missing, where they are in the network, and how far ahead one predicts. We call this displacement the \emph{restoration shift}: the difference between the forecast the same model would make from the complete history and the forecast it makes from the observed part. If the restoration shift could be estimated at deployment time, it could simply be added back.

Estimating the restoration shift is hard because, by definition, the complete history of the query is not available. This paper argues that two kinds of evidence are nevertheless available and that they are complementary.
\emph{Historical evidence} is plentiful: a traffic network repeats its daily and weekly patterns, so complete historical windows that resemble the observed part of the query are informative about how the query will evolve, and, crucially, one can artificially impose the query's own failure pattern on those windows and observe how the forecaster reacts.
\emph{Parametric evidence} is available in training: for training windows both the incomplete and the complete history exist, so a model can be taught to predict the restoration shift directly from the incomplete input.
We expect the value of each source to depend on how much of the neighbourhood is still observed. When most of the neighbourhood reports, the observed part of the query identifies good historical analogues and their future evolution transfers well. When most of the neighbourhood is dark, analogues become ambiguous, but the forecaster's reaction to the failure pattern itself, which is what replay and counterfactual restoration capture, becomes the dominant source of error.

Building on this view we propose \textbf{ReShift}, a correction framework placed on top of a missing-aware STLAformer backbone. Its contributions are:
\begin{enumerate}[leftmargin=*,label=(\arabic*)]
\item We formalise the \emph{restoration shift} of a forecaster under an observation mask and use a frozen copy of the pretrained forecaster as a teacher that defines it on training data (Section~\ref{sec:problem}).
\item We introduce \emph{evolution retrieval}, which transfers anchored future evolutions of retrieved complete training windows; retrieval keys of incomplete queries are aligned with those of complete windows by a contrastive objective with a momentum-updated memory encoder (Section~\ref{sec:evolution}).
\item We introduce \emph{failure replay}, which imposes the query's own mask on complete training windows that resemble its observed part and transfers their teacher restoration shifts, i.e.\ empirical evidence about how this particular failure pattern displaces the forecast (Section~\ref{sec:replay}).
\item We add a \emph{counterfactual restoration} head that predicts the query's own restoration shift and is distilled from the teacher on complete training histories (Section~\ref{sec:cf}).
\item We combine the three corrections with \emph{coverage-conditioned allocation gates} that see only label-free observability descriptors, after per-sensor, per-horizon reranking against an explicit no-correction option and risk gating (Section~\ref{sec:allocation}).
\item We evaluate under a leakage-safe protocol with five failure families and missing rates from 20\% to 90\%, independent ablations and in-checkpoint interventions that test whether each component uses its evidence as intended (Sections~\ref{sec:experiments}--\ref{sec:results}).
\end{enumerate}

The remainder of the paper is organised as follows. Section~\ref{sec:related} reviews related work. Section~\ref{sec:problem} defines the problem and the restoration shift. Section~\ref{sec:method} presents ReShift. Section~\ref{sec:experiments} describes the experimental protocol, Section~\ref{sec:results} reports results, Section~\ref{sec:discussion} discusses limitations and Section~\ref{sec:conclusion} concludes.
